\usepackage{tikz}
\usetikzlibrary{positioning}
\usetikzlibrary{arrows}
\usetikzlibrary{calc} %%
%https://tex.stackexchange.com/questions/23536/how-can-you-position-a-node-relatively-to-another-in-tikz
\usetikzlibrary {shapes.multipart}
\usetikzlibrary {shapes.geometric}
\usetikzlibrary{patterns}
\usetikzlibrary{decorations.markings}
\usetikzlibrary{decorations.pathreplacing}
\usetikzlibrary{calligraphy}
\usetikzlibrary{overlay-beamer-styles}
\usetikzlibrary{fadings}
\usepackage{circuitikz}

%\tikzstyle{window} = [rectangle, rounded corners, text centered, draw=black, fill=white, minimum width = 10cm, minimum height = 8cm, anchor=south west]

%\tikzstyle{other} = [rectangle, rounded corners, text centered, draw=black, fill=white, anchor=south west, minimum width=#1 cm, minimum height=#2 cm]

% Scaling factor
%\newcommand\scalingFactorTikz{0.45}

% To resize TikZ picture with resizebox
\newcommand\resizeTB[2]{\scalebox{#1\linewidth}{!}{#2}}
\newcommand\resizeTikz[1]{\resizebox{0.45\linewidth}{!}{#1}}

% Create the window
\newcommand\drawWindow{
    %\draw[draw=white, fill=black!10] (0,9) [sharp corners] -- (0,10) [rounded corners] -- (14,10) [rounded corners] -- (14,9) [sharp corners];
    \node[window] at (0,0) {};
    \node[entry={10cm}{0.6cm}] at (0.25,9.5) {};
    \draw (0,9) -- (14,9);
}

% For TikZ styles
\tikzset{
% for every node
every text node part/.style={align=center},
every node/.style={outer sep=0pt}, % avoid anchor being a bit shifted
% text-only node
text_only/.style={anchor=west},
% rounded block node
r_block/.style 2 args={rectangle, rounded corners, text centered, draw=black, fill=white, minimum width = #1, minimum height = #2, anchor=west},
% sharp block node
s_block/.style 2 args={rectangle, text centered, draw=black, fill=white, minimum width = #1, minimum height = #2, anchor=west},
memory/.style={rectangle split, rectangle split parts=4, rectangle split part align=center, draw=black, thin, text width=2.5cm, text centered, execute at begin node=\strut, outer sep=2pt}
}

\tikzset{
    make origin horizontal center of bounding box/.style={%
        execute at end picture={%
            \path
                let 
                    \p1=(current bounding box.west),
                    \p2=(current bounding box.east)
                    in ({-max(-1*\x1,\x2)},\y1) ({max(-1*\x1,\x2)},\y1);
        }
    }
}

\newcommand{\scopedrawsequence}[9][]{
    \begin{scope}
        \foreach \idx/\val [count=\yi from 0]
            in {0/#2, 1/#3, 2/#4, 3/#5, 4/#6, 5/#7, 6/#8, 7/#9}
        {
            \node (\idx) at (#1,-\yi) {\val};
        }
        \draw[->]
            (0) edge (1)
            (1) edge (2)
            (2) edge (3)
            (3) edge (4)
            (4) edge (5)
            (5) edge (6)
            (6) edge (7);
    \end{scope}
}
